\directlua{start=os.clock();texio.write_nl("=====>Hello from Lua ! ")}
\documentclass{article}
% ne pas utiliser standalone avec les primitives pdf 
\usepackage[paperwidth=20cm, paperheight=20cm,margin=0cm]{geometry}
\pagestyle{empty}
\usepackage{tikz}
\usepackage{luacode}
\begin{document}
\noindent\begin{tikzpicture}
\begin{luacode*}
checkpoint = os.clock()
texio.write_nl("\n=====> Début du bloc Lua au bout de  " .. (checkpoint - start) .. "s.")

local N = 4000
local scale = 0.005
local PIXEL_SIZE = 0.02
local squares = {}
local max = math.floor(math.sqrt(2* N*N))
for i = 1,max  do
  squares[i*i] = true
end

local function draw_point(x,y)
	return string.format(
        "\\fill[black] (%.3f,%.3f) rectangle (%.3f,%.3f);",
        x,y,x+PIXEL_SIZE,y+PIXEL_SIZE
      )
end

local linesToPrint = {}

for x = 1, N do
  for y = 1, x do
    local s = x*x + y*y
    if squares[s] then
    	table.insert(linesToPrint,draw_point(x*scale,y*scale))
    	table.insert(linesToPrint,draw_point(y*scale,x*scale))
    end
  end
end

texio.write_nl("\n=====> Triplets calculés et insérés dans la table 'linesToPrint' en " .. (os.clock() - checkpoint) .. "s.")
checkpoint = os.clock()

tex.print(table.concat(linesToPrint," "))
\end{luacode*}
\end{tikzpicture}

\directlua{texio.write_nl("=====> Arrivée à la fin du document en : " .. (os.clock() - checkpoint) .. "s.");checkpoint = os.clock()}
\directlua{texio.write_nl("=====> Temps écoulé : " .. (os.clock() - start) .. "s.")}
\end{document}